\documentclass[11pt]{article}
\usepackage[T1]{fontenc}
\usepackage{lmodern}
\usepackage[margin=1in]{geometry}
\usepackage{amsmath,amssymb,amsthm}
\usepackage{booktabs}
\usepackage[hidelinks]{hyperref}
\usepackage{microtype}
\newtheorem{theorem}{Theorem}[section]
\newtheorem{lemma}[theorem]{Lemma}
\newtheorem{proposition}[theorem]{Proposition}
\newtheorem{corollary}[theorem]{Corollary}
\theoremstyle{definition}
\newtheorem{definition}[theorem]{Definition}
\theoremstyle{remark}
\newtheorem{remark}[theorem]{Remark}
\newcommand{\N}{\mathbb{N}}
\newcommand{\Z}{\mathbb{Z}}
\newcommand{\Ap}{\operatorname{Ap}}
\newcommand{\Max}{\operatorname{Max}}
\newcommand{\PC}{\operatorname{PC}}
\newcommand{\ind}[1]{\mathbf{1}_{\{#1\}}}
\newcommand{\pos}[1]{(#1)_{+}}

\title{A path certificate for Wilf's conjecture\\and verification up to multiplicity 23}
\author{Jason Collier}
\date{5 October 2026}

\begin{document}
\maketitle

\begin{abstract}
We give a lower bound for the Wilf number of a numerical semigroup in
terms of a finite downset of Ap\'ery normal forms and the residues of its
minimal generators. The bound is a weighted longest-path certificate
and is uniform in the generator heights. A discrete summation identity
and a carry estimate yield the shape criterion $W(S)\ge R_0(B)-D(B)$.
Two independent exhaustive enumerations, each rerun in full,
verify this criterion on every required cyclic tile of multiplicities
$20$ through $23$, with one exception at multiplicity $21$. That exception
is certified directly at all maximal targets in its unit orbit.
Combining these calculations with established cases proves Wilf's
conjecture for all numerical semigroups of multiplicity at most $23$,
without a bound on genus or conductor.
\end{abstract}

\section{Introduction}

A numerical semigroup is a submonoid $S\subseteq\N=\{0,1,2,\ldots\}$
with finite complement. For $S\ne\N$, write $m$ for its multiplicity,
$e$ for the cardinality of its minimal generating set, and $c$ for its
conductor. Put $n=|S\cap[0,c)|$. Wilf's conjecture asks whether
\begin{equation}\label{eq:wilf}
 W(S):=en-c\ge0
\end{equation}
always holds \cite{Wilf}.

Bruns, Garc\'ia-S\'anchez, O'Neill and Wilburne established the
conjecture for $m\le18$ using a polyhedral algorithm \cite{BGOW}.
Kliem and Stump extended this to $m=19$ \cite{KS}. The cases $e\le3$
are classical \cite{FGH,DM}; Eliahou proved the conjecture whenever
$3e\ge m$ \cite{EliahouGraph}. Genus-based computation has verified it
through genus $100$ \cite{DEF}. The August 2026 account \cite{Marashdeh}
still lists $m\le19$ as the fixed-multiplicity bound.

Here we use finite downsets of exponent vectors representing the
Ap\'ery set. Downsets, also viewed as sets of monomials closed under
division, already play a role in the literature; see, for example,
\cite{EliahouDivsets}. Our purpose is to give a self-contained path
certificate on such a downset and apply it to a finite census. The
construction is uniform in the integer heights of the generators, so a
finite calculation addresses infinitely many semigroups.

\begin{theorem}[Computer-assisted]\label{thm:main}
Every numerical semigroup of multiplicity $m\le23$ satisfies Wilf's
conjecture.
\end{theorem}

For $20\le m\le23$, the only cases not supplied by $e\le3$ and
$3e\ge m$ are
\[
\begin{array}{c|c}
m&\text{embedding dimensions requiring computation}\\ \hline
20,21&4,5,6\\
22,23&4,5,6,7.
\end{array}
\]
These are the fourteen rows of Table~\ref{tab:census}. No assertion for
unrestricted multiplicity follows from this finite computation.

\section{Ap\'ery normal forms and a summation identity}

Let
\[
S=\langle m,g_1,\ldots,g_d\rangle,\qquad e=d+1,
\]
be a minimal generating set. Each $g_j>m$ has the form
$g_j=a_j+mh_j$, with $1\le a_j<m$ and integer $h_j\ge1$.
The $a_j$ are distinct: two generators in the same residue class would
make the larger one decomposable. Also $\gcd(m,a_1,\ldots,a_d)=1$.

The Ap\'ery set
\[
\Ap(S,m)=\{s\in S:s-m\notin S\}
\]
contains one element in each residue class modulo $m$. For each of its
elements, select the lexicographically least factorization over
$g_1,\ldots,g_d$. Let $B\subseteq\N^d$ be the set of selected exponent
vectors, and put $w(b)=\sum_jg_jb_j$.

\begin{lemma}\label{lem:normal}
The set $B$ is a downset of cardinality $m$, contains the origin and all
unit vectors, and the map
$\phi(b)=\sum_ja_jb_j\pmod m$ is bijective on $B$.
\end{lemma}
\begin{proof}
If $b\in B$ and $v\le b$ coordinatewise, then $w(v)\in\Ap(S,m)$:
otherwise $w(b)-m=(w(v)-m)+w(b-v)$ would belong to $S$.
If a factorization of $w(v)$ were lexicographically smaller than $v$,
adding $b-v$ would give a smaller factorization of $w(b)$.
Thus $v\in B$. Each $g_j$ lies in $\Ap(S,m)$, since otherwise
$g_j=m+(g_j-m)$ would be decomposable, and by minimality $\mathbf e_j$
is its only factorization over $g_1,\ldots,g_d$. Hence $B$ contains the
unit vectors, and the origin, the factorization of $0$. Distinct
elements of $\Ap(S,m)$ have distinct selected factorizations, so
$|B|=m$. Finally $w(b)\equiv\phi(b)\pmod m$, and $w$ maps $B$ onto
$\Ap(S,m)$, which meets every residue class modulo $m$ exactly once.
Therefore $\phi$ is bijective on $B$.
\end{proof}

\begin{definition}
A \emph{nondegenerate cyclic tile} of size $m$ in $\N^d$ is a downset
$B$ of cardinality $m$ containing the origin and all unit vectors,
together with labels $a\in\{1,\ldots,m-1\}^d$ such that
$\phi(b)=a\cdot b\pmod m$ is bijective on $B$.
\end{definition}

The census includes every such tile, whether or not it arises from a
numerical semigroup. This larger class suffices by Lemma~\ref{lem:normal}.
Write $|b|=\sum_jb_j$ and define the endpoint weight
\begin{equation}\label{eq:mu}
\mu(b)=\sum_{j:b+\mathbf e_j\notin B}(b_j+1).
\end{equation}
In particular, $\mu(b)\ge0$.

\begin{lemma}[Discrete summation]\label{lem:summation}
For every function $f:B\to\mathbb R$,
\begin{equation}\label{eq:summation}
\sum_{b\in B}\mu(b)f(b)
=d\sum_{b\in B}f(b)
+\sum_{b\in B}\sum_{j:b_j>0}b_j\bigl(f(b)-f(b-\mathbf e_j)\bigr).
\end{equation}
Consequently
\begin{equation}\label{eq:moments}
\sum_b\mu(b)=dm,\qquad
\sum_b\mu(b)b_j=(d+1)\sum_bb_j\quad(1\le j\le d).
\end{equation}
\end{lemma}
\begin{proof}
Expand the left side of \eqref{eq:summation} as
\[
\sum_b(|b|+d)f(b)
-\sum_{\substack{b,j\\b+\mathbf e_j\in B}}(b_j+1)f(b).
\]
In the second sum, substitute $v=b+\mathbf e_j$. The downset property
identifies the indexing set with all $(v,j)$ satisfying $v_j>0$.
This gives \eqref{eq:summation}. Apply it to $f=1$ and $f(b)=b_j$.
\end{proof}

\section{The path certificate}

For a cyclic tile, put $A(b)=a\cdot b$. At a target
$t\in\{0,\ldots,m-1\}$, set
\begin{align}
l_t&=m-1-t,& N_t(b)&=\left\lfloor\frac{A(b)+l_t}{m}\right\rfloor,
\label{eq:carry}\\
\Psi_t(b)&=|b|+N_t(b),&
E_t(b)&=\max_{z\in B,\,z\ge b}\bigl(\Psi_t(z)-\Psi_t(b)\bigr).
\label{eq:path}
\end{align}
The maximum in \eqref{eq:path} is attained at a coordinatewise maximal
cell: $\Psi_t$ increases strictly at each upward unit step. Define
\begin{align}
p_t&=l_t+\sum_b\mu(b)N_t(b)-(d+1)\sum_bN_t(b),\label{eq:p}\\
R_t&=\sum_b\mu(b)E_t(b),& \PC_t&=p_t+R_t.\label{eq:pc}
\end{align}
All these quantities are integers computed from the finite labelled tile.

For a genuine semigroup, let $\beta=\max\Ap(S,m)=c+m-1$ and let $y\in B$
be its selected factorization. The cell $y$ is maximal, since the
generators are positive. Write $r=\beta\pmod m$, $q=\lfloor\beta/m\rfloor$,
and
\[
\ell_b=\left\lfloor\frac{\beta-w(b)}m\right\rfloor\ge0.
\]

\begin{proposition}[Exact endpoint identity]\label{prop:identity}
With the notation above,
\begin{equation}\label{eq:identity}
W(S)=p_r+\sum_b\mu(b)\ell_b.
\end{equation}
\end{proposition}
\begin{proof}
There are exactly $\ell_b$ members of
$w(b)+m\N$ below $c=\beta-m+1$. Thus $n=\sum_b\ell_b$.
Since $c=mq-l_r$ and $w(b)=A(b)+m h\cdot b$,
\[
\ell_b=q-h\cdot b-N_r(b).
\]
It follows that
\[
W=dmq-(d+1)\sum_b h\cdot b-(d+1)\sum_bN_r(b)+l_r.
\]
The two identities in \eqref{eq:moments} give
\[
\sum_b\mu(b)\ell_b
=dmq-(d+1)\sum_b h\cdot b-\sum_b\mu(b)N_r(b).
\]
Subtract and use \eqref{eq:p}.
\end{proof}

\begin{theorem}[Path bound]\label{thm:path}
For every numerical semigroup with the normal-form tile above,
\begin{equation}\label{eq:pathbound}
\ell_b\ge E_r(b)\quad(b\in B),\qquad W(S)\ge\PC_r.
\end{equation}
\end{theorem}
\begin{proof}
Let $z\in B$ with $z\ge b$. The formula for $\ell$ in the proof of
Proposition~\ref{prop:identity} and $h_j\ge1$ give
\[
\ell_b-\ell_z
=h\cdot(z-b)+N_r(z)-N_r(b)
\ge |z|-|b|+N_r(z)-N_r(b).
\]
As $\ell_z\ge0$, maximizing over $z$ proves the first inequality.
Multiply it by the nonnegative weights $\mu(b)$ and apply
Proposition~\ref{prop:identity}.
\end{proof}

The path bound uses only $h_j\ge1$ and $\ell_z\ge0$; it requires no
upper bound on the heights.

\section{A criterion depending only on the shape}

Define
\begin{equation}\label{eq:shape}
\eta(b)=\max_{z\in B,\,z\ge b}(|z|-|b|),\quad
R_0(B)=\sum_b\mu(b)\eta(b),\quad
D(B)=\sum_b\pos{|b|-2}.
\end{equation}
Here $x_+=\max(x,0)$; the subscript in $R_0$ indicates a shape quantity,
not the target $t=0$ in $R_t$.

\begin{lemma}[Carry estimate]\label{lem:carry}
For every cyclic tile and every target $t$, one has $p_t\ge-D(B)$.
\end{lemma}
\begin{proof}
Put $\rho_t(b)=(A(b)+l_t)\pmod m$ and
$s_t(b)=\sum_{j:a_j>\rho_t(b)}b_j$.
If $b_j>0$, subtracting $a_j$ shows that
\[
N_t(b)-N_t(b-\mathbf e_j)=\ind{a_j>\rho_t(b)}.
\]
By \eqref{eq:summation},
$p_t=l_t-\sum_b(N_t(b)-s_t(b))$.
Since $\rho_t$ is a permutation of the residues, precisely $l_t$ of its
values exceed $t$. Therefore
\begin{equation}\label{eq:slot}
p_t=-\sum_b F_t(b),\qquad
F_t(b)=N_t(b)-s_t(b)-\ind{\rho_t(b)>t}.
\end{equation}

We claim $F_t(b)\le |b|-2$ when $|b|\ge2$, and $F_t(b)\le0$ when
$|b|\le1$. For $b=0$ this is immediate. For a unit vector,
$N_t(b)\in\{0,1\}$; if it is $1$, then $\rho_t(b)<a_j$ and
$s_t(b)=1$, proving the claim.

Now write $k=|b|\ge2$, $s=s_t(b)$, $\rho=\rho_t(b)$, $N=N_t(b)$,
and $l=l_t$. If $s=k$, then
$Nm+\rho=A(b)+l\le k(m-1)+l<(k+1)m$, so $N\le k$ and $F_t(b)\le0$.
If $s<k$, the $k-s$ noncarrying label occurrences are at most $\rho$,
and the remaining $s$ are at most $m-1$. Hence
\[
Nm\le(k-s-1)\rho+s(m-1)+l
\le(k-1)(m-1)+l<km.
\]
Thus $N\le k-1$. If $s\ge1$, or if $N\le k-2$, the claim follows.
The remaining possibility is $s=0$ and $N=k-1$. In that case
$(k-1)(m-\rho)\le l$. If $\rho\le t$, the left side is at least
$(k-1)(m-t)>m-1-t=l$, a contradiction. Consequently $\rho>t$, and
$F_t(b)=k-2$. Summing the claim in \eqref{eq:slot} proves the lemma.
\end{proof}

\begin{lemma}\label{lem:height}
For every cyclic tile and target, $R_t\ge R_0(B)$.
\end{lemma}
\begin{proof}
The labels are positive, so $N_t(z)\ge N_t(b)$ for $z\ge b$.
It follows that $E_t(b)\ge\eta(b)$. Sum with weights $\mu(b)$.
\end{proof}

\begin{corollary}[Shape certificate]\label{cor:shape}
Every numerical semigroup satisfies
\begin{equation}\label{eq:shapebound}
W(S)\ge\PC_r\ge R_0(B)-D(B).
\end{equation}
In particular, $R_0(B)\ge D(B)$ suffices for Wilf's conjecture, for
every realization of the labelled tile and every genus.
\end{corollary}

The criterion is sufficient; no universal assertion $R_0(B)\ge D(B)$
is made. The next section includes a cyclic tile where that inequality
fails and the path certificate still succeeds.

\section{The finite verification}

\subsection{Finite candidate sets and exhaustive enumeration}

For any nondegenerate downset of size $m$, every nonzero $b\in B$ obeys
\begin{equation}\label{eq:candidates}
|b|\le m-d,\qquad
\prod_{j=1}^d(b_j+1)+\#\{j:b_j=0\}\le m.
\end{equation}
Indeed, a monotone path to $b$ contains $|b|+1$ cells and exactly one
unit vector, so adjoining the other $d-1$ units gives the first bound.
The box $[0,b]$ has the displayed product as its cardinality; the units
in zero coordinates lie outside that box, giving the second bound.
The origin is included separately. These inequalities give an explicit
finite candidate set containing every possible tile.

Both production programs enumerate strictly increasing label vectors
with entries in $\{1,\ldots,m-1\}$ and with
$\gcd(m,a_1,\ldots,a_d)=1$. For a unit
$u\in(\Z/m\Z)^\times$, multiplying labels by $u$ preserves bijectivity.
Sorting the resulting labels is a simultaneous coordinate permutation
of the tile. One lexicographically least sorted label vector is chosen
from each unit orbit. The orbit weight is the number of distinct sorted
label vectors in that orbit, which may be smaller than $\varphi(m)$.
The shape quantities $R_0$ and $D$ are invariant under coordinate
permutation. These reductions therefore preserve all required shapes.

The first enumerator, archived as \texttt{tiles\_enum.py}, constructs
every cell whose box has injective residue labels. Such cells form a
downset, and each has box volume at most $m$, so the construction is
finite. Every cell of a cyclic tile must occur in this pool.
The search assigns one cell to each residue. Choosing a candidate forces
all cells in its box; a branch is rejected only when a residue would be
assigned two different cells or an unassigned residue has no consistent
candidate. It chooses the next residue deterministically with the fewest
consistent candidates. For any fixed valid tile, its assigned cell is
always available along the branch corresponding to that tile, so the
search finds every tile. The deterministic branching and forced
assignments also give a unique branch for each completed tile. Finally
the nondegeneracy condition is imposed. The production calls have no
tile cap or degree cutoff.

The second enumerator uses \eqref{eq:candidates}, orders candidate cells
by degree and then lexicographically, and extends the forced origin and
units by an increasing sequence. A cell may be added only when its
immediate parents have already been selected and its residue is unused.
Every downset has exactly one such increasing sequence. A branch is
pruned when an unused residue has no later candidate, or every later
candidate of that residue has an already excluded parent. Both tests
are necessary conditions for an extension and cannot discard a tile.
This search shares no exact-cover enumeration code with the first one.

\subsection{Completed census rows}

Table~\ref{tab:census} contains exactly the cases needed beyond the
established theorems. The weighted count is the sum of label-orbit
weights, not the number of representatives actually visited. In fresh
full runs, the two programs give identical weighted counts, minima of
$R_0-D$ and exception counts on every row; the representative counts
are from the first program, as the second does not report them.
The representative count includes every tile at the chosen label
vectors; it does not take a further quotient by their stabilizers.

\begin{table}[htbp]
\centering
\small
\begin{tabular}{rrrrrr}
\toprule
$m$ & $d=e-1$ & Representatives & Weighted tiles & $\min(R_0-D)$ & Exceptions\\
\midrule
20&3&1,391&10,244&4&0\\
20&4&25,619&187,000&12&0\\
20&5&328,112&2,492,114&14&0\\
21&3&1,383&15,846&$-2$&1\\
21&4&30,133&331,341&16&0\\
21&5&444,166&5,195,104&11&0\\
22&3&1,763&17,265&0&0\\
22&4&41,611&397,790&14&0\\
22&5&664,133&6,574,135&12&0\\
22&6&8,169,530&81,030,042&19&0\\
23&3&1,487&32,714&1&0\\
23&4&37,699&795,993&16&0\\
23&5&670,311&14,746,842&18&0\\
23&6&9,460,136&206,425,967&16&0\\
\midrule
\multicolumn{2}{l}{Total}&19,877,474&318,252,397&&1\\
\bottomrule
\end{tabular}
\caption{The fourteen required census rows. Exceptions count
representatives with $R_0<D$.}\label{tab:census}
\end{table}

Write $H=\sum_b\pos{|\operatorname{supp}(b)|-1}$. An earlier version of
the first census evaluated only $R_0-\min(H,D)$. Nonnegative values of that margin alone do not establish the criterion
used here. Both programs now record $R_0-D$ directly; these minima and
the direct treatment below supply precisely the evidence needed for
Corollary~\ref{cor:shape}. The present proof does not require a bound
involving $H$.

\subsection{The exception at multiplicity 21}

Both programs report the same sole exceptional representative in
Table~\ref{tab:census}. It has $a=(1,4,5)$ and
\begin{equation}\label{eq:exception}
\begin{split}
B={}&\{(0,j,k):j,k\ge0,\ j+k\le4\}\\
 &\cup\{(i,0,k):1\le i\le3,\ k\ge0,\ i+k\le3\}.
\end{split}
\end{equation}
It has $|B|=21$, $R_0(B)=15$ and $D(B)=17$. The eight maximal cells are
\[
\{(0,j,4-j):0\le j\le4\}\cup\{(i,0,3-i):1\le i\le3\}.
\]
For each of the twelve units $u\pmod{21}$, replace $a$ by
$a^{(u)}=(u,4u,5u)\pmod{21}$, with residues in $\{1,\ldots,20\}$,
and evaluate \eqref{eq:pc} at each maximal target
$t=a^{(u)}\cdot z\pmod{21}$. The minima over the eight targets are
\[
\begin{array}{c|rrrrrrrrrrrr}
u&1&2&4&5&8&10&11&13&16&17&19&20\\ \hline
\min_t\PC_t&15&22&31&23&29&23&23&22&25&22&23&23.
\end{array}
\]
Thus all $96$ certificates are nonnegative. Simultaneous coordinate
permutation changes none of their values. Every genuine realization of
this orbit has a maximal Ap\'ery target, so Theorem~\ref{thm:path}
certifies it.

\begin{proof}[Proof of Theorem~\ref{thm:main}]
The semigroup $\N$ is immediate.
The case $m\le19$ follows from \cite{BGOW,KS}. For $20\le m\le23$,
apply the established results when $e\le3$ or $3e\ge m$.
In every remaining case Lemma~\ref{lem:normal} gives a tile included in
Table~\ref{tab:census}. Except for \eqref{eq:exception}, it satisfies
$R_0\ge D$, so Corollary~\ref{cor:shape} applies. The preceding
calculation certifies the exceptional orbit. This exhausts all cases.
\end{proof}

\section{Reproducibility and scope}

The computational supplement, supplied as ancillary files with the
arXiv source submission, contains the two production algorithms, the
outputs of their full reruns, the earlier archived row summaries and
terminal logs, a SHA-256 manifest, and \texttt{verify.py}. Its default
checks recompute the extrema saved in $99$ archived records, compare
the fourteen required rows of the archived and rerun outputs, check
that both reruns report the same single exceptional tile, evaluate the
$96$ exceptional certificates, and verify the endpoint and path
identities on eight retained genuine semigroups and fifty seeded
examples. Its \texttt{-{}-smoke} option also performs fourteen fresh
small enumerations through $m=9$, compares the complete tile sets of
an exact-cover and a linear-extension search for each label vector,
and tests the carry estimate at every target on those tiles.

The independent program was compiled with GCC~13.3.0
(\texttt{-O3 -std=gnu11}) and run under Ubuntu~24.04 in the Windows
Subsystem for Linux; the exact-cover census was run with CPython~3.14.5
under Windows~11. The supplement records the commands, source hashes,
run times and outputs. All fourteen C runs recorded exit code zero.
Python exit codes were not captured; each run produced its final summary
record, written after enumeration, and left standard error empty.
These reruns, rather than the archived
logs, are the evidence for Table~\ref{tab:census}. The archived outputs
agree with them on every row, but some were produced by source
revisions not retained in the supplement.

Two further checks share no code with the production programs.
\texttt{shape\_first\_d3.py} enumerates every downset of size $m$ in
$\N^3$ containing the units, counts its valid labellings directly, and
reproduces the four rows with $d=3$, including the exception and its
certificates. \texttt{genus\_sweep.py} checks
Lemmas~\ref{lem:normal}, \ref{lem:summation}, \ref{lem:carry}
and~\ref{lem:height}, Proposition~\ref{prop:identity} and
Theorem~\ref{thm:path} on all $258{,}582$ numerical semigroups of genus
at most $22$, testing the carry and height bounds at every target.
Fourteen of these semigroups have $R_0(B)<D(B)$; the path certificate
covers each of them. Instructions for reproducing the computations and
checking the saved evidence are included in the supplement.

The theorem is uniform in genus and generator height. The certificate
itself applies in every dimension, but the finite census establishes
its required sign only in the stated range. Extending that range or
proving positivity for general cyclic tiles requires further work.

\paragraph{Computer and AI assistance.}
This is a computer-assisted proof: its finite verification uses the
exhaustive enumerations described above. AI tools also assisted with
mathematical exploration, proof development, code preparation and
manuscript drafting and review. The retained sources and outputs make
the computational evidence available for inspection and reproduction.
The author is responsible for the mathematical claims and the submitted
manuscript.

\begin{thebibliography}{9}
\bibitem{Wilf}
H. S. Wilf, A circle-of-lights algorithm for the ``money-changing
problem'', \emph{Amer. Math. Monthly} \textbf{85} (1978), 562--565.

\bibitem{FGH}
R. Fr\"oberg, C. Gottlieb and R. H\"aggkvist, On numerical semigroups,
\emph{Semigroup Forum} \textbf{35} (1987), 63--83.

\bibitem{DM}
D. E. Dobbs and G. L. Matthews, On a question of Wilf concerning
numerical semigroups, in: \emph{Focus on Commutative Rings Research},
Nova Sci. Publ., New York, 2006, 193--202.

\bibitem{BGOW}
W. Bruns, P. A. Garc\'ia-S\'anchez, C. O'Neill and D. Wilburne,
Wilf's conjecture in fixed multiplicity,
\emph{Internat. J. Algebra Comput.} \textbf{30} (2020), 861--882.
\href{https://arxiv.org/abs/1903.04342}{arXiv:1903.04342}.

\bibitem{KS}
J. Kliem and C. Stump, A new face iterator for polyhedra and for more
general finite locally branched lattices,
\emph{Discrete Comput. Geom.} \textbf{67} (2022), 1147--1173.
\href{https://arxiv.org/abs/1905.01945}{arXiv:1905.01945}.

\bibitem{EliahouGraph}
S. Eliahou, A graph-theoretic approach to Wilf's conjecture,
\emph{Electron. J. Combin.} \textbf{27} (2020), Paper No. 2.15.
\href{https://arxiv.org/abs/1909.03699}{arXiv:1909.03699}.

\bibitem{EliahouDivsets}
S. Eliahou, Divsets, numerical semigroups and Wilf's conjecture,
\emph{Comm. Algebra} \textbf{53} (2025), 2025--2048.

\bibitem{DEF}
M. Delgado, S. Eliahou and J. Fromentin, A verification of Wilf's
conjecture up to genus 100, \emph{J. Algebra} \textbf{664} (2025), 150--163.

\bibitem{Marashdeh}
M. F. Marashdeh, An upper bound for the type of a numerical semigroup,
and a reduction of Wilf's conjecture, preprint, August 2026.
\href{https://arxiv.org/abs/2608.12531}{arXiv:2608.12531}.
\end{thebibliography}
\end{document}
